% Options for packages loaded elsewhere
\PassOptionsToPackage{unicode}{hyperref}
\PassOptionsToPackage{hyphens}{url}
%
\documentclass[
  a4paper,
  13pt]{extarticle}
\usepackage{amsmath,amssymb}
\usepackage{setspace}
\usepackage{iftex}
\ifPDFTeX
  \usepackage[T1]{fontenc}
  \usepackage[utf8]{inputenc}
  \usepackage{textcomp} % provide euro and other symbols
\else % if luatex or xetex
  \usepackage{unicode-math} % this also loads fontspec
  \defaultfontfeatures{Scale=MatchLowercase}
  \defaultfontfeatures[\rmfamily]{Ligatures=TeX,Scale=1}
\fi
\usepackage{lmodern}
\ifPDFTeX\else
  % xetex/luatex font selection
  \setmainfont[]{Times New Roman}
\fi
% Use upquote if available, for straight quotes in verbatim environments
\IfFileExists{upquote.sty}{\usepackage{upquote}}{}
\IfFileExists{microtype.sty}{% use microtype if available
  \usepackage[]{microtype}
  \UseMicrotypeSet[protrusion]{basicmath} % disable protrusion for tt fonts
}{}
\usepackage{xcolor}
\usepackage[left=3cm,right=2cm,top=2cm,bottom=2cm]{geometry}
\setlength{\emergencystretch}{3em} % prevent overfull lines
\providecommand{\tightlist}{%
  \setlength{\itemsep}{0pt}\setlength{\parskip}{0pt}}
\setcounter{secnumdepth}{-\maxdimen} % remove section numbering
\ifLuaTeX
\usepackage[bidi=basic]{babel}
\else
\usepackage[bidi=default]{babel}
\fi
\babelprovide[main,import]{vietnamese}
\ifPDFTeX
\else
\babelfont{rm}[]{Times New Roman}
\fi
% get rid of language-specific shorthands (see #6817):
\let\LanguageShortHands\languageshorthands
\def\languageshorthands#1{}
% ctu_xelatex_header.tex — header-includes cho build LaTeX CTU (xelatex).
% Ánh xạ các ký tự Unicode mà TeX Gyre Termes thiếu glyph (chỉ số dưới/trên,
% ký tự kẻ hộp) sang lệnh LaTeX tương đương để không bị rớt glyph khi biên dịch.
\usepackage{newunicodechar}
\usepackage{pifont}
\usepackage{amssymb}
\newunicodechar{✓}{\ding{51}}
\newunicodechar{✗}{\ding{55}}
\newunicodechar{≥}{\ensuremath{\geq}}
\newunicodechar{≤}{\ensuremath{\leq}}
\newunicodechar{≈}{\ensuremath{\approx}}
\newunicodechar{↔}{\ensuremath{\leftrightarrow}}
\newunicodechar{₀}{\textsubscript{0}}
\newunicodechar{₁}{\textsubscript{1}}
\newunicodechar{₂}{\textsubscript{2}}
\newunicodechar{₃}{\textsubscript{3}}
\newunicodechar{₄}{\textsubscript{4}}
\newunicodechar{₅}{\textsubscript{5}}
\newunicodechar{₆}{\textsubscript{6}}
\newunicodechar{₇}{\textsubscript{7}}
\newunicodechar{₈}{\textsubscript{8}}
\newunicodechar{₉}{\textsubscript{9}}
\newunicodechar{²}{\textsuperscript{2}}
\newunicodechar{³}{\textsuperscript{3}}
\newunicodechar{─}{\rule[0.5ex]{0.9em}{0.4pt}}

% --- Bảng rộng: thu nhỏ cỡ chữ + khoảng cột, cho phép ngắt token dài để
%     bảng không tràn lề phải (overfull hbox). Áp dụng cho mọi longtable
%     (pandoc sinh toàn bộ bảng dạng longtable).
\usepackage{etoolbox}
\usepackage{seqsplit}      % ngắt được token rất dài (mã biến, số dài)
% Lưu ý: KHÔNG đổi \tabcolsep — pandoc tính bề rộng cột theo công thức
% (\columnwidth - 4\tabcolsep)*frac (tham chiếu \tabcolsep trực tiếp), nên
% giảm \tabcolsep lại làm cột RỘNG hơn và tràn nặng hơn.
\AtBeginEnvironment{longtable}{%
  \footnotesize
  \emergencystretch=3em
  \hbadness=10000          % không cảnh báo underfull do RaggedRight
}
% Cho phép TeX bẻ dòng linh hoạt hơn trong ô p{} hẹp
\setlength{\emergencystretch}{2em}

% URL/DOI dài: cho phép ngắt ở bất kỳ ký tự nào (tránh tràn lề ở danh mục TLTK).
% Kết hợp với pandoc +autolink_bare_uris để URL trần được bọc \url{}.
\usepackage{xurl}
\ifLuaTeX
  \usepackage{selnolig}  % disable illegal ligatures
\fi
\IfFileExists{bookmark.sty}{\usepackage{bookmark}}{\usepackage{hyperref}}
\IfFileExists{xurl.sty}{\usepackage{xurl}}{} % add URL line breaks if available
\urlstyle{same}
\hypersetup{
  pdflang={vi},
  hidelinks,
  pdfcreator={LaTeX via pandoc}}

\author{}
\date{}

\begin{document}

\setstretch{1.2}
\hypertarget{thuxf4ng-tin-vux1ec1-nhux1eefng-ux111uxf3ng-guxf3p-mux1edbi-vux1ec1-mux1eb7t-hux1ecdc-thuux1eadt-luxfd-luux1eadn-cux1ee7a-luux1eadn-uxe1n}{%
\section*{THÔNG TIN VỀ NHỮNG ĐÓNG GÓP MỚI VỀ MẶT HỌC THUẬT, LÝ LUẬN CỦA
LUẬN
ÁN}\label{thuxf4ng-tin-vux1ec1-nhux1eefng-ux111uxf3ng-guxf3p-mux1edbi-vux1ec1-mux1eb7t-hux1ecdc-thuux1eadt-luxfd-luux1eadn-cux1ee7a-luux1eadn-uxe1n}}
\addcontentsline{toc}{section}{THÔNG TIN VỀ NHỮNG ĐÓNG GÓP MỚI VỀ MẶT
HỌC THUẬT, LÝ LUẬN CỦA LUẬN ÁN}

\textbf{Tên luận án:} Quốc tế hóa và hiệu quả hoạt động kinh doanh của
các doanh nghiệp ở châu Á

\textbf{Chuyên ngành:} Quản trị kinh doanh --- \textbf{Mã số:} 9340101

\textbf{Nghiên cứu sinh:} Đỗ Thùy Hương

\textbf{Người hướng dẫn khoa học:} PGS.TS. Phan Anh Tú

\textbf{Cơ sở đào tạo:} Trường Đại học Cần Thơ

\hypertarget{nhux1eefng-ux111uxf3ng-guxf3p-mux1edbi-cux1ee7a-luux1eadn-uxe1n}{%
\subsection*{Những đóng góp mới của luận
án}\label{nhux1eefng-ux111uxf3ng-guxf3p-mux1edbi-cux1ee7a-luux1eadn-uxe1n}}
\addcontentsline{toc}{subsection}{Những đóng góp mới của luận án}

\begin{enumerate}
\def\labelenumi{\arabic{enumi}.}
\item
  \textbf{Khung lý thuyết tích hợp đa tầng (CDCM).} Luận án đề xuất
  khung \emph{Điều tiết theo Quốc gia -- Số hóa -- Năng lực}
  (Country--Digital--Capability Moderation), kết hợp bốn lý thuyết kinh
  điển (mô hình Uppsala, lý thuyết dựa trên nguồn lực, lý thuyết thể
  chế, lý thuyết thượng tầng quản trị) với lăng kính năng lực số thành
  một kiến trúc lý thuyết thống nhất và có khả năng kiểm định. Đây là
  khung tích hợp đầy đủ ba tầng điều tiết đầu tiên được xây dựng có hệ
  thống cho bối cảnh châu Á.
\item
  \textbf{Phân tách hai cấu trúc năng lực số.} Luận án tách chỉ số năng
  lực công nghệ (TCI) và chỉ số chấp nhận số (DAI) thành hai cấu trúc
  không trùng lặp và kiểm định chúng riêng biệt trong cùng một mô hình,
  làm rõ TCI là \emph{năng lực nền tảng} nâng mặt bằng đồng nhất, còn
  DAI là \emph{nguồn lực tình huống} chỉ phát huy trong những bối cảnh
  nhất định.
\item
  \textbf{Hệ phân loại thể chế ICRV sáu chế độ.} Luận án phát triển và
  áp dụng khung phân loại \emph{Biến thiên Chế độ Bối cảnh Thể chế}
  (ICRV) với sáu chế độ con cho 50 nền kinh tế châu Á và Thái Bình
  Dương, bao gồm bảy nền kinh tế đảo nhỏ Thái Bình Dương như trường hợp
  ranh giới cực đoan --- mức độ phân giải thể chế chưa từng có trong
  nghiên cứu thực nghiệm về quan hệ quốc tế hóa và hiệu quả tại châu Á.
\item
  \textbf{Cơ sở dữ liệu đa quốc gia lớn nhất ghi nhận được cho khu vực.}
  Mẫu gộp gồm 88.869 doanh nghiệp từ 50 nền kinh tế châu Á và Thái Bình
  Dương (mẫu hồi quy chính N = 81.022), mở rộng đáng kể so với các nền
  tảng dữ liệu đa quốc gia hiện có cho quan hệ quốc tế hóa và hiệu quả
  tại châu Á.
\item
  \textbf{Bằng chứng về tính ổn định theo thời gian của điểm uốn.} Luận
  án kiểm định tính bền vững của cấu trúc phi tuyến tại Trung Quốc và
  cho thấy điểm uốn trong quan hệ quốc tế hóa và hiệu quả tương đối ổn
  định qua hơn một thập kỷ chuyển đổi số (≈ 48,78\% FSTS), bất chấp
  những thay đổi lớn của bối cảnh kinh tế.
\item
  \textbf{Đề xuất và định danh khái niệm Gánh nặng quốc tế hóa bắt buộc
  (FIP).} Luận án đề xuất một điều kiện biên lý thuyết của mô hình chữ U
  ngược, trong đó quan hệ quốc tế hóa và hiệu quả chuyển sang âm đơn
  điệu khi ba điều kiện cấu trúc bị vi phạm đồng thời --- một cấu trúc
  chưa được định danh trong tài liệu kinh doanh quốc tế.
\end{enumerate}

\textbf{Tuyên bố tính mới trung tâm:} Luận án là công trình đầu tiên
trong tài liệu kinh doanh quốc tế chuyển luận điểm ``quan hệ quốc tế hóa
và hiệu quả phụ thuộc bối cảnh'' từ một nhận định định tính thành một
cấu trúc kiểm định được trên toàn bộ phổ thể chế của một khu vực, đồng
thời xác định và định danh điểm mà tại đó dạng hàm chữ U ngược không còn
ứng dụng được. Khác với các nghiên cứu trước thường dừng ở việc chứng
minh thể chế điều tiết \emph{độ lớn} của hiệu ứng, luận án chứng minh
thể chế điều tiết cả \emph{dấu} và \emph{dạng hình} của quan hệ, qua đó
nâng luận điểm từ ``phụ thuộc bối cảnh'' lên ``có điều kiện theo bối
cảnh''.

\hypertarget{uxfd-nghux129a-thux1ef1c-tiux1ec5n}{%
\subsection*{Ý nghĩa thực
tiễn}\label{uxfd-nghux129a-thux1ef1c-tiux1ec5n}}
\addcontentsline{toc}{subsection}{Ý nghĩa thực tiễn}

Bản đồ điều tiết theo chế độ thể chế cung cấp cho doanh nghiệp và nhà
hoạch định chính sách một cơ sở để định vị kỳ vọng về quan hệ quốc tế
hóa và hiệu quả theo từng bối cảnh thể chế cụ thể, thay vì áp dụng một
kỳ vọng phổ quát; đồng thời cảnh báo rằng vai trò ``lá chắn số'' của DAI
phụ thuộc bối cảnh, và chỉ ra điều kiện biên FIP đối với các nền kinh tế
đảo nhỏ.

\vspace{1.4cm}

\noindent\begin{minipage}[t]{0.48\textwidth}\centering \textbf{NGƯỜI HƯỚNG DẪN KHOA HỌC}\\[2.2cm]PGS.TS. Phan Anh Tú\end{minipage}\hfill\begin{minipage}[t]{0.48\textwidth}\centering \textbf{NGHIÊN CỨU SINH}\\[2.2cm]Đỗ Thùy Hương\end{minipage}

\newpage

\hypertarget{information-on-the-new-academic-and-theoretical-contributions-of-the-dissertation}{%
\section*{INFORMATION ON THE NEW ACADEMIC AND THEORETICAL CONTRIBUTIONS
OF THE
DISSERTATION}\label{information-on-the-new-academic-and-theoretical-contributions-of-the-dissertation}}
\addcontentsline{toc}{section}{INFORMATION ON THE NEW ACADEMIC AND
THEORETICAL CONTRIBUTIONS OF THE DISSERTATION}

\textbf{Dissertation title:} Internationalization and Firm Performance
in Asia: The Moderating Roles of Institutions, Digital Capability, and
Top-Manager Characteristics

\textbf{Major:} Business Administration --- \textbf{Code:} 9340101

\textbf{PhD candidate:} Do Thuy Huong

\textbf{Supervisor:} Assoc. Prof.~Dr.~Phan Anh Tu

\textbf{Institution:} Can Tho University

\hypertarget{new-contributions-of-the-dissertation}{%
\subsection*{New contributions of the
dissertation}\label{new-contributions-of-the-dissertation}}
\addcontentsline{toc}{subsection}{New contributions of the dissertation}

\begin{enumerate}
\def\labelenumi{\arabic{enumi}.}
\item
  \textbf{An integrative, multi-tier theoretical framework (CDCM).} The
  dissertation proposes the \emph{Country--Digital--Capability
  Moderation} framework, joining four established theories (the Uppsala
  model, the resource-based view, institutional theory, and
  upper-echelons theory) with a digital-capability extension into a
  unified, testable architecture. It is the first systematically
  constructed three-tier moderation framework for the Asian context.
\item
  \textbf{Separation of two digital-capability constructs.} The study
  separates the Technological Capability Index (TCI) and the Digital
  Adoption Index (DAI) into two non-overlapping constructs and tests
  them jointly within a single model, clarifying that TCI is a
  \emph{foundational capability} that shifts the performance level
  uniformly, whereas DAI is a \emph{situational resource} effective only
  in specific contexts.
\item
  \textbf{A six-regime institutional taxonomy (ICRV).} The dissertation
  develops and applies the \emph{Institutional Context Regime Variation}
  taxonomy of six sub-regimes for 50 Asia-Pacific economies, including
  seven Pacific Small Island Developing States as an extreme boundary
  case --- an institutional resolution not previously used in empirical
  I--P research on Asia.
\item
  \textbf{The largest documented multi-country dataset for the region.}
  The pooled sample comprises 88,869 firms from 50 Asia-Pacific
  economies (main regression sample N = 81,022), substantially extending
  existing multi-country evidence bases for the
  internationalization--performance relationship in Asia.
\item
  \textbf{Evidence on the temporal stability of the turning point.}
  Testing the durability of the nonlinear structure in China, the
  dissertation shows that the inverted-U turning point is relatively
  stable across more than a decade of digital transformation (≈ 48.78\%
  FSTS), despite substantial changes in the economic context.
\item
  \textbf{Proposal and naming of the Forced Internationalization Penalty
  (FIP).} The dissertation introduces a theoretical boundary condition
  of the inverted-U model, under which the I--P relationship turns
  monotonically negative when three structural conditions are violated
  simultaneously --- a construct not previously named in the
  international-business literature.
\end{enumerate}

\textbf{Central novelty statement:} This dissertation is the first work
in the international-business literature to convert the claim that ``the
I--P relationship is context-dependent'' from a qualitative assertion
into a construct testable across the full institutional spectrum of a
region, while identifying and naming the point at which the inverted-U
functional form ceases to apply. Unlike prior work that demonstrates
institutions moderating the \emph{magnitude} of the effect, this
dissertation shows that institutions moderate both the \emph{sign} and
the \emph{shape} of the relationship, thereby elevating the claim from
``context-dependent'' to ``context-contingent.''

\hypertarget{practical-implications}{%
\subsection*{Practical implications}\label{practical-implications}}
\addcontentsline{toc}{subsection}{Practical implications}

The regime-based moderation map provides firms and policymakers with a
basis for calibrating their expectations about the
internationalization--performance relationship to each specific
institutional context, rather than applying a universal expectation; it
also cautions that DAI's ``digital shield'' role is context-dependent
and identifies the FIP boundary condition for small-island economies.

\vspace{1.4cm}

\noindent\begin{minipage}[t]{0.48\textwidth}\centering \textbf{SUPERVISOR}\\[2.2cm]Assoc. Prof. Dr. Phan Anh Tu\end{minipage}\hfill\begin{minipage}[t]{0.48\textwidth}\centering \textbf{PhD CANDIDATE}\\[2.2cm]Do Thuy Huong\end{minipage}

\end{document}
